\scanpage{397}\mspage{311}

que je vais expliciter un peu. On a $a\Rightarrow\deleted{\unsure{b}}a'$ puisque \deleted{$\underline S$} si a est satisfaite, $\underline S$ \uline{et} $M_\pi$ sont stables \inserted{\uline{dans} $M$} par $\underset{\longrightarrow}{\lim}$ grdes devant $\pi'_0$ et de cardinal $\leq\pi$, donc aussi leur intersection, $\underline S(\pi)$, qui est donc aussi stable dans $M_\pi$ par lesdites limites. \deleted{\unsure{L|a}} L'implication $\deleted{\unsure{d}\Rightarrow\unsure{c}},$ $c\Rightarrow c'$ \deleted{\unsure{est donc}} est une tautologie, et $c\Rightarrow b$ aussi, compte tenu de 4.13.12 qui nous dit que le syst. inductif canonique $(x_i)_{i\in I}$, avec $I=M_\pi/x$, de limite $x$, est $\pi_0$-$\pi$-adapté

(NB Comme $\pi^{\pi'_0}=\pi$, \inserted{et $\pi'_0\geq\pi_0$} on a aussi $\pi^{\pi_0}=\pi$, donc $\pi$ est utile pour $M$, \deleted{\unsure{)}} de sorte que 4.13.12 s'applique.

Reste à prouver la ``$\pm$ tautologie'' $\deleted{\unsure{b}}a'\Rightarrow\deleted{\unsure{d}}c$. \deleted{\illegible{1}} \inserted{Soit} donc une représentation $\pi_0$-$\pi$-adaptée

\[
x\simeq\underset{i\in I}{\underset{\longrightarrow}{\lim}}\,x_i,
\]

prouvons que

\[
x\in\mathrm{Ob}\,\underline S\Longleftrightarrow I_{\underline S}\ \text{cofinal dans }I.
\]

% slanted
\margin{On n'utilise, \inserted{ici, pour $\Leftarrow$,} de l'hypothèse $(x_i)_{i\in I}$ $\pi_0$-$\pi$-adapté, que le seul fait que $I$ soit grd devant $\pi$.}

L'implication $\Leftarrow\deleted{\unsure{>}}$ provient du fait que, $\underline S$ étant $\pi$-accessible dans $M$, est stable \inserted{dans $M$} par $\underset{\longrightarrow}{\lim}$ grdes devant $\pi$*. Reste à prouver

% archive sheet number: 396
